
Large language models exhibit cross-correlation ($\rho$ = 0.80--0.87) across trustworthiness benchmarks, a structure often dismissed as a scale confound. This study investigates whether a shared latent factor persists after controlling for model size and training recency. Through meta-analysis of 4,561 models from the Open LLM Leaderboard, principal component analysis on residualized benchmark scores yields a dominant first component ($\lambda _1$ = 2.277, explaining 60\% of residual variance) that exceeds permutation null thresholds (95th percentile = 0.803, p = 0.001). This residual factor correlates with a Behavioral Stability Index proxy ($\rho$ = 0.405, 95\% CI [0.380, 0.429], p < $10^{-179})$. Analysis of 16 matched base/instruct model pairs shows instruction-tuning increases both BSI (Cohen's d = 1.87) and PC1 scores (d = 1.99). Prospective validity testing indicates 5 of 6 holdout benchmarks load $\geq$ 0.3 on frozen factor weights. These findings are consistent with---but do not definitively establish---a shared representational stability mechanism underlying cross-benchmark correlation. Key limitations include synthetic BSI scores and simulated holdout benchmarks; validation with real behavioral measurements is required before strong mechanistic claims can be made.
